% Research notes for integration into the parent article.
% Requires amsmath, amssymb, amsthm and theorem/lemma/proposition/corollary.
% Classical reference keys suggested: HatcherNotes, PerelmanExtinction.

\section{A terminal criterion from freeness and Euler characteristic}
\label{sec:free-euler-terminal}

The preceding terminal proposal assumed irreducibility of every whole-cut
component.  The source validator does not establish that hypothesis.  A
different, exact gate removes the need to establish irreducibility and
also removes the need for a separate boundary-component census.  The
topological ingredients below are classical.  The contribution to the
current programme is their use with a source-authenticated whole-cut
presentation and an independently authenticated Euler characteristic.

All manifolds in this section are compact orientable PL three-manifolds.
Surfaces and disks are locally flat.  Homology without specified
coefficients has rational coefficients.  Every component of the manifold
under consideration is assumed to have nonempty boundary.  This last
condition holds for the entire normal cut of the current connected
compact source manifold with nonempty boundary, even if the normal
surface has closed or one-sided components.  A complementary component
must meet the original boundary or a side of the cutting surface; a
component meeting neither would be a closed-and-open component of the
original connected manifold.

\begin{lemma}[Free group and the boundary rank]
\label{lem:free-boundary-rank}
Let $N$ be connected, with connected boundary of genus $g$.  If
$\pi_1(N)\cong F_g$, then $N$ is a genus-$g$ handlebody.
\end{lemma}
\begin{proof}
This proof makes its Poincar\'e-theorem input explicit and does not assume
irreducibility.  A positive-genus boundary surface cannot have its
fundamental group embedded in a free group: subgroups of free groups are
free, whereas closed orientable positive-genus surface groups are not.
The Loop Theorem therefore supplies a properly embedded compression
disk~\cite[Theorem 3.1]{HatcherNotes}.
Cut along such a disk whenever a positive-genus boundary component
remains.  Van Kampen's theorem expresses the original group as the free
product of the cut-component groups, together with an additional free
generator when a disk does not separate the manifold.  Thus every new
component again has free fundamental group.

The disk-cut process terminates because
\[
 \mu=\sum_{F\subset\partial N,\ g(F)>0}(2g(F)-1)
\]
strictly decreases under every compression.  An essential separating
compression decreases this quantity by one; a nonseparating compression
decreases it by two, except on a torus, where it decreases it by one.
The Loop Theorem is used only as an existence theorem in this argument;
the proposed terminal algorithm does not search for these disks.

Suppose that after $d$ disk cuts there are $v$ connected components,
with a total of $b$ spherical boundary components and free-group ranks
$r_1,\ldots,r_v$.  All these components have nonempty boundary, so
$b\geq v$.  Every disk compression adds two to the Euler characteristic
of the boundary, hence
\[
 2b=2-2g+2d,\qquad d=b+g-1.
\]
Reconstructing $N$ attaches $d$ one-handles to the $v$ pieces.  The graph
of these attachments is connected.  Van Kampen and abelianization give
\[
 g=\sum_{i=1}^v r_i+d-v+1
  =g+\sum_{i=1}^v r_i+b-v.
\]
Every term in $\sum_i r_i+b-v$ is nonnegative.  Consequently all $r_i$
are zero and $b=v$.  Each terminal component is simply connected and has
one sphere boundary.  Capping this boundary by a ball gives a closed
simply connected three-manifold.  The three-dimensional Poincar\'e
theorem~\cite{PerelmanExtinction} identifies it with $S^3$; the PL Schoenflies theorem, equivalently
Alexander's theorem in this dimension, then identifies the uncapped
component with a ball~\cite[Theorem 1.1]{HatcherNotes}.  Reversing the disk cuts constructs $N$ from balls
and one-handles, so it is a handlebody.  Its genus is $g$.
\end{proof}

\begin{theorem}[Aggregate free-group--Euler criterion]
\label{thm:free-euler}
Let
\[
 W=\coprod_{i=1}^{c}N_i
\]
be a finite disjoint union of connected compact orientable PL
three-manifolds with nonempty boundary.  Form the algebraic free product
of the component groups.  Suppose
\[
 *_{i=1}^c\pi_1(N_i)\cong F_r.
\]
Then the following are equivalent:
\begin{enumerate}
\item $\chi(W)=c-r$;
\item $H_2(W;\mathbb Q)=0$;
\item every $N_i$ is a handlebody.
\end{enumerate}
When they hold, if $\pi_1(N_i)\cong F_{r_i}$, the genus of $N_i$ is
$r_i$ and $\sum_i r_i=r$.  In particular, $W$ is a union of balls if and
only if, in addition, $r=0$.
\end{theorem}
\begin{proof}
Each free-product factor injects into $F_r$.  By the Nielsen--Schreier
theorem and finite generation, it is a finite-rank free group $F_{r_i}$.
Abelianizing the free product gives $\sum_i r_i=r$.  Since every component
has nonempty boundary, $H_3(N_i)=0$, and therefore
\begin{equation}
 \chi(W)=c-r+\sum_{i=1}^c b_2(N_i).
 \label{eq:aggregate-beta-two}
\end{equation}
This proves the equivalence of the first two assertions.

Assume $b_2(N_i)=0$, and write $b_i=|\pi_0(\partial N_i)|\geq1$.
The homology sequence of the pair contains
\[
 H_3(N_i,\partial N_i)\longrightarrow H_2(\partial N_i)
       \longrightarrow H_2(N_i)=0.
\]
The first group is one-dimensional, generated by the relative orientation
class, and the second has dimension $b_i$.  The first map has rank one:
it sends the orientation class to the sum of the boundary orientation
classes.  Exactness forces $b_i=1$.  Doubling $N_i$ along its boundary
gives a closed odd-dimensional manifold, of Euler characteristic zero;
thus $\chi(\partial N_i)=2\chi(N_i)$.  If its connected boundary has
genus $g_i$, this identity and $b_2(N_i)=0$ give
\[
 1-g_i=\chi(N_i)=1-r_i,
\]
so $g_i=r_i$.  Lemma~\ref{lem:free-boundary-rank} shows that $N_i$ is a
handlebody.  The converse follows from the homology of a handlebody.
\end{proof}

There is no need to match two independently produced component
histograms in this theorem.  It suffices to authenticate the entire
free product, the total component count, and the total Euler
characteristic.  In fact, a group certificate is allowed to mix
generators from different algebraic components.  If a genus histogram is
wanted without a separate boundary census, the implementation should
retain the source component of each generator and certify freeness
componentwise.  That stronger bookkeeping is useful for output, but is
not required for the all-handlebody conclusion.

The distinction between group topology and a geometric move is essential.
Deleting a freely trivial relator is a Tietze transformation, even though
deleting the corresponding two-handle can change the manifold.  The
Euler gate above accounts for the topology that the fundamental group
does not see.  The checker must use the source Euler characteristic, not
the Euler characteristic of a simplified presentation complex.

\paragraph{Sharpness of the gate.}
The manifold $S^2\times I$ has trivial fundamental group, one connected
component and Euler characteristic two.  It fails $\chi=c-r$, and has
two boundary components.  Similarly,
$H_g\mathbin{\#}(S^1\times S^2)$ has connected genus-$g$ boundary and
free fundamental group of rank $g+1$.  Its Euler characteristic is
$1-g$, so its excess over $c-r=-g$ is one.  These are genuine
non-handlebody obstructions invisible to freeness alone.  Conversely,
Euler characteristic and abelianization alone cannot replace the
free-group certificate; a punctured nontrivial homology sphere is the
standard obstruction.

\section{A finite whole-cut group compiler for disk parallelity bases}
\label{sec:disk-group-compiler}

Assume the existing source-parallelity and ordered-interface verifiers
have authenticated a normal cut of a $t$-tetrahedron source, and every
parallelity base is a disk.  The latter assertion includes orientability,
product structure, Euler characteristic one and exactly one boundary
component.  A ported base has multiplicity one by source authentication.
Unported disk bundles may occur with binary-encoded multiplicity.  The
construction below retains the entire source core, including all its
two- and three-handles.

Write $h_j$ for the number of source core $j$-handles, and $\kappa$ for
the number of components of its zero/one-handle graph.  Fix a spanning
forest of this graph.  Every nonforest one-handle becomes a generator.
The cyclic attaching curve of every source two-handle becomes a signed
relator by erasing tree crossings.  Signs are read from source face
flags.  Comparing only zero-handle endpoint labels is insufficient when
a one-handle is a loop; its two source face flags still distinguish its
two directions.  Attaching a three-handle along a sphere does not change
the fundamental group, so it contributes no relator and is retained in
the Euler ledger.

There are two possibilities for a ported disk bundle.  Its vertical
boundary can be an entire attaching annulus, in which case the ball
$D^2\times I$ is exactly a two-handle attached along $\partial D^2\times I$.
Trace the authenticated annular midpoint through the core and append its
signed word as a relator.  Alternatively its intersection with the core
is a nonempty disjoint union of attaching disks, corresponding to proper
interval components of its base boundary.  The bundle itself is a ball.
It therefore contributes a contractible vertex space, and each attaching
disk contributes an edge with trivial edge group in the graph-of-spaces
description.  Because a disk base has only one boundary component, a
full circle and proper interval attachments cannot occur on that same
base.  An unported disk bundle is an isolated ball.

Let $a_\circ$ be the number of full annular disk bundles, $d$ the number
of disk bundles with proper interval attachments, and $j$ the number of
those interval attachments.  Form a graph $\Gamma$ with the $\kappa$
core components and these $d$ disk balls as vertices, and the $j$
attaching disks as edges.  If $c_f=|\pi_0(\Gamma)|$, its cycle rank is
\[
 \beta=j-\kappa-d+c_f.
\]
Van Kampen gives the free product of the augmented core groups and an
additional free factor of total rank $\beta$, with the distribution
among components specified by $\Gamma$.  These additional generators
have no occurrences in any relator.  This graph calculation uses
contractible attaching disks, so no unspecified geometric slide or
primitive-curve claim is involved.

\begin{theorem}[Disk-base whole-cut presentation]
\label{thm:disk-compiler}
Under the authenticated source hypotheses just stated, the construction
produces a presentation of the algebraic free product of the fundamental
groups of all whole-cut components.  Its numbers of generators and
relators are
\begin{equation}
 a=h_1-h_0+\kappa+\beta,
 \qquad m=h_2+a_\circ.
 \label{eq:compiler-counts}
\end{equation}
If there are $b$ unported disk balls, counted with their binary
multiplicities, the number of whole-cut components is $c=c_f+b$.
Furthermore
\begin{equation}
 \chi(W)=c-a+m-h_3.
 \label{eq:source-presentation-ledger}
\end{equation}
For the current native source format there is a uniform linear bound
on $a$, $m$ and total literal signed relator length.  The conservative
bounds
\[
 a\leq74t,\qquad m\leq78t,\qquad
 \sum_{\rho}|\rho|\leq78t
\]
suffice.  No term depends on the represented normal-surface weight.
\end{theorem}
\begin{proof}
A zero/one-handle core is a regular neighborhood of its graph.  Collapsing
a spanning forest gives its standard free generators, and the usual
handle presentation theorem attaches each two-handle relation using its
signed curve.  Source flags and the annular-midpoint reconstruction
authenticate the precise words.  Three-handles do not change the group.
The preceding graph-of-spaces argument then proves the presentation claim
and~\eqref{eq:compiler-counts}, including the trivial groups of the
isolated balls.

Euler characteristic is additive across these gluings.  A full annular
disk bundle contributes $\chi(D^2\times I)-\chi(S^1\times I)=1$.
A disk bundle with proper interval attachments contributes one, with one
subtracted for every attaching disk.  Thus
\[
 \chi(W)=h_0-h_1+h_2-h_3+a_\circ+d+b-j.
\]
Substituting $\beta=j-\kappa-d+c_f$, $c=c_f+b$ and the counts of
\eqref{eq:compiler-counts} proves~\eqref{eq:source-presentation-ledger}.

The existing constructor bounds
\[
 h_0\leq6t,\ h_1\leq8t,\ h_2\leq12t,\ h_3\leq4t,
 \qquad p\leq66t,
\]
where $p$ is the number of core ports.  The sum of the lengths of all
interior edge-link wedge lists is at most $6t$, because each local
tetrahedral edge occurs in one such list.  Every list is used at most
twice, at the two endpoint ranks.  The sum of one-handle crossings in
all source two-handle attaching cycles is therefore at most $12t$.
The ordered-interface verifier partitions the ports into boundary words,
each visiting each port once; consequently the additional full annular
words have total length at most $p\leq66t$.  Erasing forest crossings
only decreases length.  There are at most $p$ full circle bundles and
at most $p$ proper interval attachments.  In particular,
$a\leq h_1+j\leq74t$ and $m\leq h_2+p\leq78t$.
\end{proof}

The active graph, relators and port incidences have $O(t)$ records, but
the compiler must also inspect every supplied base-profile record,
including unported disk families.  Let $f$ denote the number of those
records and let $E$ denote the exact total bit mass of the supplied
profile and finite-source data.  A conservative structural bound for
the delivered forest/ordering implementation is $O(t^2+f)$ graph and
record steps, together with its exact arithmetic cost on the integers
in $E$.  In particular, a minimum-label disjoint-set implementation is
not silently assigned an inverse-Ackermann or linear bound.  Source
preparation and source certificate checking have their separate costs.
If the normal coordinates have at most $B$ bits each, the finite chart
and source-label description has $O(t(B+\log(t+2)))$ bits; the finite
group presentation itself needs only $O(t\log(t+2))$ bits after generators
are indexed.  These statements do not bound the list of profiles or an
entire existing orbit certificate by the size of the chart description.

\section{Exact polynomial terminal reduction with controlled word length}
\label{sec:no-growth-tietze}

For this finite group presentation, free and cyclic word reduction may
use sign cancellations that the earlier unsigned raw model deliberately
excluded.  We use only explicit signed words.  A current relator $\rho$
admits a defining-generator elimination at $x$ when $x$ or $x^{-1}$
occurs exactly once in $\rho$, counting both signs.  Rotate its cyclic
word to
\[
 \rho=x^{\varepsilon}u,\qquad\varepsilon\in\{1,-1\},
\]
where $u$ contains no occurrence of $x$.  The relation is the definition
$x=u^{-\varepsilon}$.  Substitute this expression for every occurrence
of $x$ in every other relator, substitute its inverse for $x^{-1}$,
delete $\rho$ and $x$, then freely and cyclically reduce the other
relators.  A freely trivial relator can be discarded; its deletion is
recorded in the count $z$ below.

The proposed deterministic producer examines available defining-generator
eliminations and accepts one only if the total cyclically reduced relator
length does not increase.  It stops when every relator is gone or no
admitted candidate remains.  All source generators, including ones absent
from every relator, remain in the generator ledger until explicitly
eliminated.  A final empty relator list presents the free group whose rank
is the number of surviving generators.

\begin{theorem}[Soundness and polynomial local complexity]
\label{thm:bounded-tietze}
Let $a,m$ be the numbers of source generators and relators, and $L$ the
total literal source word length.  Use compact integer generator and row
labels, or charge the input-label dictionary separately.  Put
$N=a+m+L+1$.  The procedure above
is a sound one-sided freeness test and terminates after at most
$\min(a,m)$ defining-generator eliminations.  A straightforward exact
implementation uses $O(N^4)$ indexed word operations for discovery,
$O(N^3)$ for independent replay of a supplied elimination trace, and
$O(N^2)$ working letter storage.  A trace recording source row labels
and eliminated generator labels has $O(m\log(N+1))$ bits; all
substitutions and trivial-relator drops can be reconstructed.
\end{theorem}
\begin{proof}
Free reduction leaves the represented relator unchanged.  Cyclic rotation
or cyclic reduction replaces it by a conjugate, so its normal closure is
unchanged.  A unique occurrence of $x$ makes $\rho$ an explicit
definition of that generator.  Substitution followed by deleting the
definition is an elementary Tietze elimination, with inverse operation
adjoining the generator and its defining relation.  Deleting a freely
trivial relator removes an identity.  These facts prove soundness,
independently of the length policy.  Each elimination removes one
generator and one relator, proving termination.

At every accepted state the total reduced relator length is at most $L$.
There are at most $L$ candidate row/generator pairs, since each contains a
distinct occurrence within its donor row.  For one candidate, the donor
tail has length at most $L$, and at most $L$ letters in all other words
are replaced.  The total temporary substituted length is at most
$O(L^2)$.  A stack reduction and endpoint trim reduce those words in
linear time in that temporary length.  Accounting also for generator and
row ledgers costs $O(N^2)$ per candidate.  Testing at most $N$ candidates
at at most $N$ elimination stages gives $O(N^4)$ discovery work.

A verifier receives the chosen row/generator pair at each stage and
needs no candidate search.  It reconstructs the donor solution, all
substitutions, reductions, empty-relator deletions, the length gate and
the final generator ledger.  Thus it uses $O(N^3)$ work.  At most
$O(N^2)$ temporary letters are required.  Stable labels use
$O(\log(N+1))$ bits, so the listed pairs have the asserted size.  Source
authentication and any appended source proof are additional to this
small algebraic trace.
\end{proof}

This is a complexity theorem for a specified one-sided producer, not a
decision theorem for group freeness.  An inconclusive result may reflect
the absence of a suitable generator, the length policy, or an
unfortunate earlier choice.  No order-independence or completeness
assertion is made for this greedy signed procedure.  The broader
certificate format may admit other exactly checkable Tietze moves, but
their discovery costs require separate bounds.

\begin{example}[A free presentation with no initial admitted move]
Define an automorphism of $F(a,b)$ by
\[
 \varphi(a)=ab,\qquad\varphi(b)=a.
\]
Its inverse is given by $a\mapsto b$ and $b\mapsto b^{-1}a$, as direct
substitution verifies.  The two words
\[
 \varphi^4(a)=abaababa,\qquad\varphi^4(b)=abaab
\]
therefore form a free basis.  In particular,
\[
 \langle a,b\mid abaababa,abaab\rangle
\]
is the trivial group: applying the inverse automorphism four times takes
its relators to $a,b$.  Both displayed words are positive and cyclically
reduced.  Their $(a,b)$ occurrence counts are $(5,3)$ and $(3,2)$, so
neither relator has a generator occurring once.  The proposed producer
has no initial defining-generator move and returns inconclusive.
Adjoining an unused generator $c$ gives the same refusal for a
presentation of $F_1$.  This is an exact algebraic counterexample to
completeness and needs no group-recognition oracle.
\end{example}

There is nevertheless a precise inclusion at the level of certificates.
Every complete unsigned column-leaf trace has a signed Tietze trace that
never increases length: its pivot generator occurs once in the entire
current relator family, so deleting its donor requires no substitution
in another relator.  Free and cyclic reductions cannot delete that unique
occurrence before its pivot.  This proves inclusion of the old
certificate class; it does not prove that an arbitrary greedy signed
choice preserves the old successful ordering.  Retaining the older
producer as a fast path makes the inclusion operational.

For a strict algebraic contrast, take
\[
 \langle x,y\mid xyx^{-1},\ yxy^{-1}\rangle.
\]
Its unsigned incidence matrix is
\[
 \begin{pmatrix}2&1\\1&2\end{pmatrix},
\]
whose two rows form a stopping set.  Cyclic reduction instead gives the
relators $y$ and $x$, and two defining-generator eliminations certify the
trivial group.  This is an algebraic illustration; the native experiment
provides the separate source-authenticated geometric example.

\begin{corollary}[Resolving the earlier linear-deferral ball family]
\label{cor:linear-deferral-resolved}
The geometric ball family of report 102~\cite{Report102}, consisting after forest
contraction of $q$ blocks with unsigned matrix
$\left(\begin{smallmatrix}1&1\\1&2\end{smallmatrix}\right)$, has a
signed terminal certificate with $2q$ defining-generator eliminations,
no trivial-relator deletions, and strictly decreasing relator length.
Its entire $2q$-row matrix is a raw stopping set, and its old raw
deferral number is $q$.  The signed certificate has
$O(q\log(q+2))$ bits.
\end{corollary}
\begin{proof}
Start with one zero-handle and two canceling one/two-handle pairs, a
handle presentation of a ball.  The initial two-handle words are $x,y$.
Slide the first two-handle over the second along the standard positive
band, and then slide the second over the new first.  Choose the bands
in the zero-handle boundary away from all cocore boundaries.  Equivalently,
perform these disk slides in the complementary handlebody of the
standard genus-two Heegaard splitting of $S^3$.  The complementary
disk system remains complete, so the resulting two-handle attachment
still gives a ball.  This is the geometric realization from report 102;
it is not inferred from words alone.

With these positive bands, the two words are $xy,yxy$.  Cyclically
rotating the second gives $xy^2$.  They are initially positive and
cyclically reduced.  Eliminate $x$ using $xy$, so
$x=y^{-1}$; the second relator freely reduces to $y$.  Eliminate $y$.
The total lengths are $5,1,0$.  There are no empty-relator deletions.
Take $q$ copies and form their boundary connected sum along disks
disjoint from all attaching regions.  The result is a ball.  After
contracting the connecting zero/one forest its matrix is block diagonal,
and its word alphabets are disjoint.  The two signed moves therefore
apply independently in every block.

Each block's two rows form a stopping set, so the whole $2q$-row set is
stopping.  Any raw deferral scheme must remove at least one row from
each of the $q$ disjoint block obstructions.  Conversely, deferring the
second row of each block leaves a singleton pivot for the first.  Hence
the raw deferral number is exactly $q$, as proved in report 102.
The $2q$ signed row/generator choices use $O(\log(q+2))$ bits per move,
proving the certificate bound.  The new conclusion is the short signed
resolution of the earlier geometric obstruction family.
\end{proof}

The family is geometric, but is not asserted to have been exported by
the native normal-cut source constructor.  The separate native target
establishes actual additional coverage in that narrower source class.
The displayed block decomposition permits a mathematically linear
blockwise construction and replay, but the delivered global solver is
only assigned the $O(N^4)$ discovery and $O(N^3)$ replay bounds above.
No linear replay implementation is claimed by the block corollary.

Abelianization is insufficient even when an exponent-sum determinant is
one.  Consider
\[
 \mathcal P=\langle x,y\mid x^3y^{-5},\ (xy)^2x^{-3}\rangle.
\]
The exponent-sum matrix is
\[
 \begin{pmatrix}3&-5\\-1&2\end{pmatrix}
\]
with determinant one.  The group is nevertheless nontrivial: under the
usual right-to-left permutation convention, assign
\[
 x=(1\ 2\ 3),\qquad y=(1\ 3\ 5\ 2\ 4).
\]
Then $x^3=y^5=1$ and $xy=(2\ 4)(3\ 5)$, so both relators vanish in
the permutation group, while $x$ has nonidentity image.  This exact
witness explains why the terminal test certifies the entire group and
does not replace it by its exponent-sum matrix.

\section{Three-handles give an exact allowance for trivial relators}
\label{sec:relator-allowance}

The source count identity~\eqref{eq:source-presentation-ledger} gives a
particularly compact interpretation of the Euler gate.  It is useful
both for implementation and for diagnosing a refusal.

\begin{theorem}[Relator-deletion ledger and early homology obstruction]
\label{thm:relator-ledger}
Start with the disk-base whole-cut presentation of
Theorem~\ref{thm:disk-compiler}.  Suppose a valid signed trace has
performed $k$ defining-generator eliminations and discarded $z$ freely
trivial relators, including all such initial normalizations.  Then
\begin{equation}
 b_2(W;\mathbb Q)\geq z-h_3.
 \label{eq:early-relator-defect}
\end{equation}
If all relators have disappeared, equality holds:
\begin{equation}
 b_2(W;\mathbb Q)=z-h_3.
 \label{eq:complete-relator-defect}
\end{equation}
In particular, a completed free-group trace passes the Euler gate
exactly when $z=h_3$.  If $z>h_3$ at any intermediate stage, the
authenticated whole cut already has positive second Betti number and
at least one component is not a handlebody.
\end{theorem}
\begin{proof}
The current presentation has
\[
 a'=a-k,\qquad m'=m-k-z,
 \qquad a'-m'=a-m+z.
\]
Its abelianization has at most $m'$ independent linear relations on $a'$
generators, so
\[
 b_1(W;\mathbb Q)\geq a'-m'=a-m+z.
\]
All components have nonempty boundary; consequently
$b_2(W)=\chi(W)-c+b_1(W)$.  Insert
$\chi(W)=c-a+m-h_3$ to obtain~\eqref{eq:early-relator-defect}.

If $m'=0$, the remaining presentation is free, of rank
$r=a'=a-k$, and $m=k+z$.  Substitution into the same Euler formula gives
\[
 b_2(W)=\chi(W)-c+r=m-k-h_3=z-h_3.
\]
The aggregate criterion of Theorem~\ref{thm:free-euler} now proves the
last conclusions.
\end{proof}

Thus when the core has no three-handles, a complete trace with no
freely-trivial relator deletions automatically satisfies the terminal
Euler condition.  A three-handle is not silently discarded from the
manifold: its contribution is recorded by the necessary equality
$z=h_3$.  A completed trace with $z<h_3$ would contradict the authenticated
source hypotheses, because it would give negative $b_2$; it is not a
legitimate alternative terminal topology.  A partial trace with
$z>h_3$ is a genuine obstruction to the whole cut being a union of
handlebodies, but has no direct implication that the original knot is
nontrivial: the chosen cut itself may simply be unsuitable for the
intended hierarchy.

Binary families cause no growth in this argument.  Adding $b$ isolated
disk balls adds $b$ to both $c$ and $\chi$, contributes no generator or
relator, and leaves every equation above invariant.  Only the bit length
of $b$, not its value, enters the arithmetic cost.

\paragraph{Terminal scope.}
The successful certificate establishes an abstract homeomorphism type
and, with componentwise group ledgers, a genus census.  It supplies no
meridian disk system in marked source coordinates and no transported
boundary pattern.  Its local polynomial cost does not provide the
hierarchy construction, a bounded marked successor, progress, or the
global depth bound required for quasi-polynomial unknot recognition.

% Source references checked in this task:
% Hatcher, Notes on Basic 3-Manifold Topology, current linked 75-page PDF:
% https://pi.math.cornell.edu/~Hatcher/3M/3Mfds.pdf
% Theorem 1.1 printed p.1 (Alexander); Theorem 3.1 printed p.56 (Loop).
% The PDF retains pre-Perelman historical text; do not cite that old status.
% Perelman, Finite extinction time for the solutions to the Ricci flow
% on certain three-manifolds, arXiv:math/0307245 (2003).
% https://arxiv.org/abs/math/0307245
% Clay official solved status and account:
% https://www.claymath.org/millennium/poincare-conjecture/
